\documentclass[12pt,a4paper]{article}
\usepackage[utf8]{inputenc}
\usepackage{amsmath}
\usepackage{amssymb}
\usepackage{amsthm}
\usepackage{graphicx}
\usepackage{tikz}
\usepackage{algorithm}
\usepackage{algpseudocode}
\usepackage{geometry}
\usepackage{hyperref}
\usepackage{cite}
\usepackage{float}


\geometry{margin=1in}

\newtheorem{theorem}{Theorem}
\newtheorem{lemma}[theorem]{Lemma}
\newtheorem{proposition}[theorem]{Proposition}
\newtheorem{corollary}[theorem]{Corollary}
\theoremstyle{definition}
\newtheorem{definition}{Definition}
\theoremstyle{remark}
\newtheorem{remark}{Remark}

\title{A Fast Algorithm for Generating All Triangulations of Biconnected Plane Graphs}
\author{Tawkir Aziz Rahman\\
Department of Computer Science and Engineering\\
Bangladesh University of Engineering and Technology (BUET)\\
Dhaka-1000, Bangladesh}
\date{\today}

\begin{document}

\maketitle

\begin{abstract}
In this paper, we deal with the problem of generating all triangulations for biconnected planar graph. We give an algorithm for generating all triangulations for a biconnected planar graph G of n vertices. Our algorithm establishes a double recursive tree structure among all triangulations of G, and generates each triangulation in O(1) amortized time using O(n) space. This is the first algorithm that achieves constant time generation for biconnected planar graphs. The key innovations of our algorithm are: (1) the concept of \emph{safe root vertex} for each face, which guarantees the existence of a root triangulation that does not create multi-edges when combined with previously triangulated faces; (2) the use of local genealogical trees for each face, allowing independent face processing while maintaining global chord consistency; and (3) a dynamic global chord set to efficiently track and prevent multi-edge conflicts during the triangulation process. We prove the correctness, completeness, and efficiency of our algorithm, and provide implementation details demonstrating its practical effectiveness.
\end{abstract}

\section{Introduction}

Triangulation of planar graphs is a fundamental problem in computational geometry with wide-ranging applications in VLSI floorplanning~\cite{kozminski1985}, computational geometry~\cite{berg2008}, and graph drawing~\cite{nishizeki2004}. While finding a single triangulation of a planar graph can be done efficiently, the problem of \emph{generating all} triangulations presents significant computational challenges. This paper addresses the problem of exhaustively generating all triangulations of biconnected planar graphs—a problem that has remained open despite substantial progress on restricted graph classes.

\subsection{Motivation and Challenges}

The enumeration of all triangulations is crucial for applications requiring exploration of the entire solution space, such as finding optimal triangulations under various cost metrics, studying structural properties of triangulation spaces, and analyzing the diameter of triangulation graphs. However, several fundamental difficulties make this problem challenging:

\textbf{Exponential solution space.} The number of triangulations for a given graph can be exponential in the number of vertices, making explicit storage and exhaustive enumeration computationally expensive. Any practical algorithm must generate triangulations on-the-fly without storing all of them simultaneously.

\textbf{Duplicate avoidance.} Naive enumeration approaches often generate the same triangulation multiple times. Traditional methods for avoiding duplicates require storing previously generated triangulations and performing isomorphism checks, which introduces substantial memory overhead (often exponential space) and adds at least logarithmic time complexity per triangulation for lookup operations.

\textbf{Efficiency requirements.} For enumeration algorithms, the gold standard is to achieve output-sensitive time complexity, ideally constant or polynomial time per generated solution. This requires careful algorithm design to avoid redundant computation.

\subsection{Prior Work and the Triconnected Case}

Classical enumeration algorithms employ various strategies to address these challenges. The generate-and-filter approach produces candidate objects and removes duplicates through filtering, but requires significant memory. Orderly generation~\cite{read1978} ensures canonical representations through sophisticated techniques but involves expensive isomorphism testing. The reverse search paradigm~\cite{avis1996} addresses memory concerns by implicitly defining a spanning tree over the solution space, enabling polynomial-space enumeration, but typically requires $O(n^2)$ time per solution for triangulation problems.

For the restricted case of \emph{convex polygons}, Hurtado and Noy~\cite{hurtado1999} constructed a graph representing all triangulations and studied its structural properties, though their approach requires generating all triangulations of smaller polygons first. For more general graphs, Li and Nakano~\cite{li2001} presented an algorithm for generating all biconnected ``based'' plane triangulations with at most $n$ vertices, but it builds incrementally from graphs with fewer vertices, making it inefficient when only triangulations of a specific $n$-vertex graph are needed.

The breakthrough for \emph{triconnected} plane graphs came from Parvez, Rahman, and Nakano~\cite{parvez2011}, who developed an elegant algorithm achieving $O(1)$ time per triangulation using $O(n)$ space. Their key insight is to organize all triangulations into a tree structure where parent-child relationships correspond to single edge flip operations. By carefully defining a unique \emph{generating chord} for each triangulation, they ensure that flipping this chord produces a unique parent, thereby avoiding duplicates without explicit storage or isomorphism testing. The algorithm decomposes the graph into faces, triangulates each face independently, and combines results efficiently.

\subsection{The Multi-Edge Problem in Biconnected Graphs}

The Parvez-Rahman-Nakano algorithm critically relies on the graph being \emph{triconnected}. Triconnectivity guarantees that any two faces share at most their boundary edges—they cannot share internal chords. This property enables independent face triangulation without conflicts.

However, for general \emph{biconnected} graphs (which have connectivity exactly 2), this guarantee fails. Biconnected graphs may contain \emph{separation pairs}—pairs of vertices whose removal disconnects the graph. Faces separated by such pairs can share internal structure beyond boundary edges. When triangulating faces independently, chords added in one face may conflict with chords needed in another face, potentially creating \textbf{multi-edges} (multiple edges between the same vertex pair), which violates the simple graph property.

This \textbf{multi-edge problem} is the fundamental obstacle preventing direct application of the triconnected algorithm to biconnected graphs. Simply applying the face-by-face triangulation approach of Parvez et al. to biconnected graphs can yield invalid triangulations with parallel edges, making the extension non-trivial.

\subsection{Our Contribution}

In this paper, we present the first algorithm that generates all triangulations of biconnected planar graphs in $O(1)$ amortized time per triangulation using $O(n)$ space. Our algorithm successfully extends the tree-of-triangulations framework to the biconnected case by introducing novel techniques to resolve the multi-edge problem:

\begin{enumerate}
    \item \textbf{Safe root vertex selection:} For each face, we introduce the concept of a \emph{safe root vertex}—a vertex from which all initial chords can be drawn without creating multi-edges with previously triangulated faces. We prove that such vertices always exist and provide an efficient method to find them.
    
    \item \textbf{Local genealogical trees:} Instead of a single global tree, we construct \emph{local genealogical trees} for each face, organizing all valid triangulations of that face. These local trees are explored via depth-first search during recursive face processing.
    
    \item \textbf{Dynamic global chord tracking:} We maintain a \emph{global chord set} that tracks all chords added across all faces during the recursive process. This enables efficient detection and prevention of multi-edge conflicts when triangulating each face.
\end{enumerate}

We rigorously prove that our algorithm guarantees:
\begin{itemize}
    \item \textbf{Completeness:} Every valid triangulation is generated exactly once (no duplicates, no omissions).
    \item \textbf{Correctness:} All generated triangulations are valid (no multi-edges, proper triangulation).
    \item \textbf{Efficiency:} $O(1)$ amortized time per triangulation with $O(n)$ space.
\end{itemize}

We also provide implementation details and experimental results demonstrating the algorithm's practical effectiveness on various graph instances.

\subsection{Paper Organization}

The remainder of this paper is organized as follows: Section 2 provides essential definitions and preliminaries on planar graphs and triangulations. Section 3 presents our algorithm in detail, including the safe root vertex selection, local genealogical tree construction, and the main recursive enumeration procedure. Section 4 provides rigorous proofs of correctness and completeness. Section 5 analyzes time and space complexity. Section 6 discusses implementation details and experimental results. Section 7 concludes with future research directions.



\subsection{Triangulation Enumeration}

Early work on triangulation enumeration focused on specific graph classes. Chazelle~\cite{chazelle1991} gave a linear-time algorithm for finding \emph{a} triangulation of a simple polygon, but not for enumerating all triangulations.

Hurtado and Noy~\cite{hurtado1999} studied the graph of triangulations of convex polygons, where vertices represent triangulations and edges represent edge flips. They established theoretical properties but did not provide efficient enumeration algorithms for arbitrary $n$-vertex polygons without first generating all smaller cases.

Li and Nakano~\cite{li2001} developed an algorithm for generating all biconnected plane triangulations with at most $n$ vertices, but this requires generating all smaller triangulations first—inefficient when only triangulations of exactly $n$ vertices are needed.

\subsection{Reverse Search and Tree-Based Enumeration}

The reverse search paradigm~\cite{avis1996} has been successfully applied to many enumeration problems. Avis presented algorithms for enumerating triangulations using reverse search, but these achieve $O(n^2)$ time per triangulation for cycles~\cite{avis1996}.

Orderly generation~\cite{read1978} avoids duplicates by outputting only canonical representations, but requires expensive isomorphism testing.

\subsection{The Parvez-Rahman-Nakano Algorithm}

The breakthrough work of Parvez, Rahman, and Nakano~\cite{parvez2011} introduced the \emph{tree of triangulations} for triconnected plane graphs, achieving $O(1)$ time per triangulation. Their key insights include:

\begin{itemize}
    \item \textbf{Tree Structure:} Triangulations form a tree where each child is generated from its parent by a single edge flip.
    \item \textbf{Root Triangulation:} The tree is rooted at a triangulation where all chords emanate from a single vertex.
    \item \textbf{Generating Chord:} Each triangulation has a unique \emph{generating chord}—the leftmost chord that can be flipped.
    \item \textbf{Parent-Child Relationship:} Flipping the generating chord produces the parent, ensuring unique relationships and no duplicates.
\end{itemize}

However, this algorithm is specifically designed for \emph{triconnected} graphs. The triconnectivity guarantee ensures that faces share only boundary edges, preventing multi-edge conflicts. For biconnected graphs without this guarantee, the algorithm fails.

\subsection{The Multi-Edge Problem}

In biconnected graphs, faces can share internal chords from previous triangulation steps. When triangulating a new face, adding a chord that already exists as a chord in a different face creates a multi-edge, violating the simple graph property. This is the fundamental barrier preventing direct application of the Parvez-Rahman-Nakano algorithm to biconnected graphs.

Our algorithm solves this problem through safe root selection and global chord tracking.

\subsection{Comparison: Triconnected vs. Biconnected Approach}

The key differences between the Parvez-Rahman-Nakano algorithm for triconnected graphs and our algorithm for biconnected graphs are:

\begin{table}
\centering
\small
\begin{tabular}{|p{3cm}|p{5cm}|p{5cm}|}
\hline
\textbf{Aspect} & \textbf{Triconnected (PRN)} & \textbf{Biconnected (Ours)} \\
\hline
Tree Structure & Single global tree of triangulations & Local trees per face + DFS combination \\
\hline
Root Selection & Any vertex can be root & Must find safe root vertex \\
\hline
Multi-edge Problem & Does not occur (guaranteed by triconnectivity) & Explicitly handled via GlobalChordSet \\
\hline
Face Processing & Can process independently & Must process sequentially with backtracking \\
\hline
Chord Conflicts & No conflicts between faces & Dynamic conflict checking required \\
\hline
Complexity & $O(1)$ per triangulation & $O(1)$ amortized per triangulation \\
\hline
Applicability & Only triconnected graphs & All biconnected graphs \\
\hline
\end{tabular}
\caption{Comparison of triconnected and biconnected triangulation algorithms}
\label{tab:comparison}
\end{table}

Our algorithm generalizes the triconnected case: when applied to a triconnected graph, every vertex is a safe root (since no chords exist between faces), and the algorithm behaves similarly to the Parvez-Rahman-Nakano algorithm.

\section{Preliminaries}
\label{sec:preliminaries}

In this section, we establish the fundamental graph-theoretic concepts and notation used throughout this paper.

\subsection{Basic Graph Theory}

Let $G = (V,E)$ denote a connected simple graph with vertex set $V$ and edge set $E$. We denote an edge connecting vertices $v_i$ and $v_j$ in $V$ by $(v_i,v_j)$. The \emph{degree} of a vertex $v$ is the number of edges incident to $v$ in $G$. 

The \emph{connectivity} $\kappa(G)$ of a graph $G$ is the minimum number of vertices whose removal results in either a disconnected graph or a single-vertex graph. A graph is \emph{$k$-connected} if $\kappa(G) \geq k$. In particular, a graph is \emph{biconnected} if it is 2-connected (i.e., it has no cut vertices) and \emph{triconnected} if it is 3-connected (i.e., it has no separation pairs).

\subsection{Planar Graphs and Plane Embeddings}

A graph $G$ is \emph{planar} if it can be embedded in the plane such that no two edges intersect geometrically except at vertices to which they are both incident. A \emph{plane graph} is a planar graph with a fixed embedding in the plane. A plane graph divides the plane into connected regions called \emph{faces}. The unbounded face is called the \emph{outer face}, and all other faces are called \emph{inner faces}. 

A plane graph is called a \emph{plane triangulation} if every face boundary contains exactly three edges. In this paper, we allow edges to be represented as non-straight curves, as is standard in topological graph theory.

\subsection{Cycles and Chords}

A \emph{cycle} $C$ in a graph $G$ is a sequence of distinct vertices $v_1,v_2,\ldots,v_k$ such that $(v_1,v_k) \in E$ and $(v_i,v_{i+1}) \in E$ for $1 \leq i < k$. In a plane graph, a cycle $C$ divides the plane into two regions: the \emph{inner region} (the region enclosed by $C$) and the \emph{outer region} (the region outside $C$). When a graph $G$ itself forms a cycle, both its inner and outer regions are faces.

Let $C$ be a cycle corresponding to the boundary of a face $F$ in a plane graph $G$, and let $v_1, v_2, \ldots, v_n$ denote the vertices on $C$ in counterclockwise order. We represent $C$ by listing its vertices as $C = \langle v_1,v_2,\ldots,v_n \rangle$. 

An edge $(v_i, v_j)$ between two vertices $v_i$ and $v_j$ on $C$ is called a \emph{chord} of $C$ if:
\begin{itemize}
    \item $(v_i,v_j) \notin E(C)$ (i.e., it is not an edge of the cycle), and
    \item $(v_i,v_j)$ is contained entirely within the face $F$ bounded by $C$.
\end{itemize}

A chord $(v_i,v_j)$ partitions the cycle $C$ into two smaller cycles: one containing vertices $v_i,v_{i+1},\ldots,v_j$ and the other containing vertices $v_j,v_{j+1},\ldots,v_i$.

\subsection{Triangulations}

A \emph{triangulation of a cycle} $C$ is a maximal set of non-crossing chords that decomposes the inner region of $C$ into triangles. In a triangulation $T$ of $C$, the chord set is maximal in the sense that any chord not in $T$ must intersect at least one chord in $T$. The sides of the resulting triangles are either chords from $T$ or edges from the original cycle $C$.

For a cycle $C$ with $n$ vertices, any triangulation contains exactly $n-3$ chords and produces exactly $n-2$ triangular faces. We represent each triangulation $T$ of $C$ by the set of its chords. Given this chord set, the triangulation can be uniquely reconstructed.

A triangulation of a cycle $C$ is called a \emph{labeled triangulation} if the vertices of $C$ are explicitly numbered sequentially from $v_1$ to $v_n$. If the vertices are not numbered, the triangulations are called \emph{unlabeled triangulations}.

\subsection{Visibility in Triangulations}

In a triangulation $T$ of a cycle $C$, we say that a vertex $y$ is \emph{visible from} a vertex $x$ if there exists a chord $(x,y)$ in $T$. This notion of visibility is fundamental to the structure of triangulations and plays a key role in our algorithm's construction of root triangulations.

\subsection{Key Definitions for Our Algorithm}

We now introduce several specialized definitions that are central to our algorithm for biconnected graphs.

\begin{definition}[Multi-edge]
A \emph{multi-edge} occurs when two or more distinct edges connect the same pair of vertices in a graph. A \emph{simple graph} contains no multi-edges or self-loops. All triangulations considered in this paper must be simple graphs.
\end{definition}

\begin{definition}[Separation Pair]
A \emph{separation pair} in a biconnected graph $G$ is an unordered pair of vertices $\{u,v\}$ whose removal disconnects $G$ into two or more components.
\end{definition}

\begin{definition}[Safe Root Vertex]
A vertex $r$ in a face $F$ of a biconnected plane graph is called a \emph{safe root vertex} if the star triangulation rooted at $r$ (where all chords emanate from $r$) does not create any multi-edge with the chords already present in previously triangulated faces.
\end{definition}

\begin{definition}[Genealogical Tree]
A \emph{genealogical tree} of triangulations for a cycle $C$ is a rooted tree structure where:
\begin{itemize}
    \item Each node represents a distinct triangulation of $C$,
    \item The root is the star triangulation from a designated root vertex,
    \item Each edge represents a single chord flip operation that transforms one triangulation into another, and
    \item The parent-child relationship is defined by a unique generating chord.
\end{itemize}
\end{definition}

\begin{definition}[Global Chord Set]
The \emph{global chord set} is a dynamic data structure (typically implemented as a hash set) that maintains all chords currently present in the partially triangulated graph during the recursive enumeration process. This set enables constant-time detection and prevention of multi-edge formation.
\end{definition}

\begin{definition}[Hash Set]
A \emph{hash set} is a data structure implementing the set abstract data type using a hash table, providing average-case $O(1)$ time complexity for insertion, deletion, and membership testing operations.
\end{definition}

These definitions provide the foundation for understanding our algorithm's approach to handling the multi-edge problem in biconnected graphs.


\section{The Algorithm}
\label{sec:algorithm}

\subsection{Algorithm Overview}

Our algorithm generates all triangulations of a biconnected plane graph $G$ with $k$ faces $F_1, F_2, \ldots, F_k$ using a recursive depth-first search (DFS) strategy. The key idea is to:

\begin{enumerate}
    \item Process faces sequentially from $F_1$ to $F_k$.
    \item For each face $F_i$, construct a \emph{local genealogical tree} of all its valid triangulations.
    \item Maintain a \emph{global chord set} to prevent multi-edge conflicts.
    \item Recursively combine face triangulations to produce complete graph triangulations.
\end{enumerate}

\subsection{Key Innovations}

Our algorithm introduces three key innovations to extend the Parvez-Rahman-Nakano framework from triconnected to biconnected graphs:

\subsubsection{Innovation 1: Local Genealogical Trees}

Instead of constructing a single global genealogical tree for the entire graph (which fails for biconnected graphs due to multi-edge conflicts), we build \emph{local genealogical trees} for each face independently.

\begin{itemize}
    \item Each face $F_i$ has its own tree $\mathcal{T}(F_i)$ of triangulations
    \item The root of $\mathcal{T}(F_i)$ is the safe root triangulation
    \item Parent-child relationships are defined by edge flips within the face
    \item Local trees are explored via DFS during recursive face processing
\end{itemize}

This decomposition enables independent triangulation of each face while maintaining global consistency through the chord set.

\subsubsection{Innovation 2: Safe Root Vertex Selection}

To ensure that the root triangulation of each face does not create multi-edges with previously triangulated faces, we introduce the concept of a \emph{safe root vertex}.

The safe root vertex is chosen such that all initial chords from this vertex to non-neighboring vertices of the face do not conflict with chords already in the global set. This guarantees a valid starting point for the local genealogical tree.

We prove (Theorem~\ref{thm:safe_root_exists}) that every face, regardless of how many previous faces have been triangulated, contains at least two safe root vertices.

\subsubsection{Innovation 3: Global Chord Set with Dynamic Management}

The \emph{global chord set} $S$ tracks all chords currently used in the partial triangulation being explored. It serves two purposes:

\begin{enumerate}
    \item \textbf{Conflict Detection:} Before adding any chord to a face triangulation, check if it already exists in $S$
    \item \textbf{Safe Root Finding:} Used to identify which vertices can serve as safe roots for the current face
\end{enumerate}

The set is managed dynamically during DFS:
\begin{itemize}
    \item \textbf{Push:} When exploring a triangulation, add its chords to $S$
    \item \textbf{Recurse:} Process next face with updated $S$
    \item \textbf{Pop:} After exploring all descendants, remove chords from $S$ (backtrack)
\end{itemize}

This ensures that $S$ always reflects exactly the chords in the current exploration path.

\subsection{Safe Root Vertex Selection}

The first critical innovation is identifying a \emph{safe root vertex} for each face.

\begin{definition}[Safe Root Vertex]
A vertex $r$ of a face $F_i$ is a \emph{safe root vertex} if adding chords from $r$ to all non-neighboring vertices of $F_i$ does not create any edge that already exists in the global chord set (from previous faces $F_1, \ldots, F_{i-1}$).
\end{definition}

\begin{algorithm}
\caption{Finding a Safe Root Vertex}
\label{alg:safe_root}
\begin{algorithmic}[1]
\Require Face $F$ with vertices $v_1, v_2, \ldots, v_n$, global chord set $S$
\Ensure A safe root vertex $r$
\For{$i = 1$ to $n$}
    \State $\text{safe} \gets \texttt{true}$
    \For{each non-neighbor $v_j$ of $v_i$ in $F$}
        \If{$(v_i, v_j) \in S$}
            \State $\text{safe} \gets \texttt{false}$
            \State \textbf{break}
        \EndIf
    \EndFor
    \If{$\text{safe}$}
        \State \Return $v_i$
    \EndIf
\EndFor
\end{algorithmic}
\end{algorithm}

We will prove in Section~\ref{sec:correctness} that every face has at least two safe root vertices.

\subsection{Root Triangulation}

Once a safe root vertex $r$ is identified, we construct the \emph{root triangulation} of the face by adding chords from $r$ to all non-neighboring vertices:

\begin{definition}[Root Triangulation]
The \emph{root triangulation} $T_r$ of a face $F$ with safe root $r$ is the triangulation formed by adding all chords $(r, v)$ for every vertex $v \in F$ that is not a neighbor of $r$ on the cycle boundary.
\end{definition}

\subsection{Local Genealogical Tree}

For each face, we build a \emph{local genealogical tree} of triangulations rooted at the safe root triangulation. The tree structure is defined by parent-child relationships based on edge flipping:

\begin{definition}[Generating Chord]
In a triangulation $T$ of a face with root vertex $r$, a chord $(v_i, v_j)$ is a \emph{generating chord} if:
\begin{enumerate}
    \item Neither $v_i$ nor $v_j$ is the root vertex $r$, and
    \item $(v_i, v_j)$ is the leftmost such chord (smallest index among non-root chords).
\end{enumerate}
\end{definition}

\begin{definition}[Parent-Child Relationship]
Let $T$ be a triangulation with generating chord $(v_i, v_j)$ forming two triangles $(v_i, v_j, v_p)$ and $(v_i, v_j, v_q)$ on opposite sides. The \emph{parent} of $T$ is the triangulation obtained by flipping $(v_i, v_j)$ to $(v_p, v_q)$, provided this flip does not create a multi-edge.
\end{definition}

\subsection{Main Algorithm}

\begin{algorithm}
\caption{Generate All Triangulations of Biconnected Plane Graph}
\label{alg:main}
\begin{algorithmic}[1]
\Require Biconnected plane graph $G$ with faces $F_1, \ldots, F_k$
\Ensure All triangulations of $G$
\State Initialize global chord set $S \gets \emptyset$
\State \Call{ProcessFace}{$1, \emptyset, S$}

\Procedure{ProcessFace}{$i, T_{\text{current}}, S$}
    \If{$i > k$}
        \State \textbf{output} $T_{\text{current}}$ \Comment{All faces triangulated}
        \State \Return
    \EndIf
    \State $r \gets$ \Call{FindSafeRoot}{$F_i, S$}
    \State $T_{\text{root}} \gets$ \Call{RootTriangulation}{$F_i, r$}
    \State \Call{Explore}{$F_i, T_{\text{root}}, i, T_{\text{current}}, S$}
\EndProcedure

\Procedure{Explore}{$F, T, i, T_{\text{current}}, S$}
    \State Add chords of $T$ to $S$
    \State \Call{ProcessFace}{$i + 1, T_{\text{current}} \cup T, S$}
    \State Remove chords of $T$ from $S$
    
    \For{each child $T'$ of $T$ in local tree}
        \If{chords of $T'$ do not conflict with $S$}
            \State \Call{Explore}{$F, T', i, T_{\text{current}}, S$}
        \EndIf
    \EndFor
\EndProcedure
\end{algorithmic}
\end{algorithm}

The algorithm performs a depth-first exploration where:
\begin{itemize}
    \item At each level $i$, we process face $F_i$
    \item For each valid triangulation $T$ of $F_i$, we recursively process $F_{i+1}$
    \item Backtracking removes chords from $S$ to explore alternative triangulations
    \item Complete triangulations are output when $i > k$
\end{itemize}

\subsection{Child Generation via Edge Flipping}

To generate children in the local genealogical tree, we flip the generating chordd:
This

\begin{algorithm}
\caption{Generate Children of Triangulation}
\label{alg:children}
\begin{algorithmic}[1]
\Require Triangulation $T$ of face $F$ with root $r$
\Ensure All children of $T$ in the local tree
\State Find generating chord $(v_i, v_j)$ of $T$
\If{no generating chord exists}
    \State \Return $\emptyset$ \Comment{$T$ is a leaf node}
\EndIf
\State Find $v_p, v_q$ forming triangles with $(v_i, v_j)$
\State $T' \gets T \setminus \{(v_i, v_j)\} \cup \{(v_p, v_q)\}$
\State \Return $\{T'\}$
\end{algorithmic}
\end{algorithm}

\subsection{Complete Algorithm Specification}

We now present the complete algorithm for generating all triangulations of a biconnected plane graph. The algorithm consists of five main procedures that work together to achieve duplicate-free enumeration.


\subsubsection{Main Entry Point}

\begin{algorithm}
\caption{Start Triangulation Process}
\label{alg:start}
\begin{algorithmic}[1]
\Procedure{StartTriangulation}{$G, k$}
    \State Label the faces $F_1, F_2, \ldots, F_k$ arbitrarily
    \State Initialize global chord set $S \gets \emptyset$
    \State $T \gets \emptyset$ \Comment{Empty partial triangulation}
    \State \Call{FindAllTriangulationsCycle}{$F, T, 0$}
\EndProcedure
\end{algorithmic}
\end{algorithm}


\subsubsection{Face-Level Triangulation}

The following procedure processes each face by finding a safe root, constructing the root triangulation, and exploring the local genealogical tree.

\begin{algorithm}
\caption{Find All Triangulations for Current Face}
\label{alg:face_triangulation}
\begin{algorithmic}[1]
\Procedure{FindAllTriangulationsCycle}{$F, T, i$}
    \State $n \gets$ number of vertices of face $F_i$
    \State $T_{ir} \gets$ \Call{FindSafeRoot}{$F, i$} \Comment{Time: $O(n)$}
    \State Add all chords of $T_{ir}$ to $S$ \Comment{Time: $O(n)$}
    \State \Call{OutputTriangulation}{$T, T_{ir}, i$}
    \State $T_i \gets T_{ir}$
    \For{$j = n - 1$ \textbf{downto} $3$}
        \State Add chord $(v_1, v_j)$ to $S$
        \State \Call{FindAllChildTriangulationsCycle}{$T_i(v_1, v_j), T, i$}
        \State Remove chord $(v_1, v_j)$ from $S$
    \EndFor
    \State Remove all chords of $T_{ir}$ from $S$ \Comment{Time: $O(n)$}
\EndProcedure
\end{algorithmic}
\end{algorithm}



\subsubsection{Recursive Child Generation}

This procedure recursively generates all children in the local genealogical tree by identifying and flipping generating chords.

\begin{algorithm}
\caption{Find All Child Triangulations in Local Tree}
\label{alg:child_triangulations}
\begin{algorithmic}[1]
\Procedure{FindAllChildTriangulationsCycle}{$T_i, T, i$}
    \State \Call{OutputTriangulation}{$T, T_i, i$}
    \If{$T_i$ has no generating chord}
        \State \Return \Comment{Leaf node reached}
    \EndIf
    \State Let $(v_b, v_{b'})$ be the leftmost generating chord of $T_i$
    \ForAll{$j \geq b$}
        \If{$(v_1, v_j)$ is a chord of $T_i$ \textbf{or} $(v_1, v_j) \notin S$}
            \State Add chord $(v_1, v_j)$ to $S$
            \State \Call{FindAllChildTriangulationsCycle}{$T_i(v_1, v_j), T, i$}
            \State Remove chord $(v_1, v_j)$ from $S$
        \EndIf
    \EndFor
\EndProcedure
\end{algorithmic}
\end{algorithm}


\subsubsection{Output and Recursion Control}

When a face is fully processed, this procedure either outputs the complete triangulation (if all faces are done) or recursively processes the next face.

\begin{algorithm}
\caption{Output Triangulation or Continue to Next Face}
\label{alg:output}
\begin{algorithmic}[1]
\Procedure{OutputTriangulation}{$T, T_i, i$}
    \If{$i = k$} \Comment{All faces triangulated}
        \State \textbf{output} $(T \cup T_i)$
    \Else
        \State \Call{FindAllTriangulationsCycle}{$F, T \cup T_i, i + 1$}
    \EndIf
\EndProcedure
\end{algorithmic}
\end{algorithm}


\subsubsection{Safe Root Finding}

The safe root finding procedure identifies a vertex from which all chords can be drawn without creating multi-edges.

\begin{algorithm}
\caption{Find Safe Root Triangulation}
\label{alg:safe_root_complete}
\begin{algorithmic}[1]
\Procedure{FindSafeRoot}{$F, \text{index}$}
    \State $r \gets$ \Call{FindSafeRootVertex}{$F, \text{index}$}
    \State $T_r \gets \emptyset$
    \State Let $F = \langle v_1, v_2, \ldots, v_n \rangle$ with $v_r = $ safe root vertex
    \For{each vertex $v_j$ in $F$}
        \If{$v_j$ is not adjacent to $v_r$ on the cycle boundary}
            \State Add chord $(v_r, v_j)$ to $T_r$
        \EndIf
    \EndFor
    \State \Return $T_r$
\EndProcedure

\Procedure{FindSafeRootVertex}{$F, \text{index}$}
    \State Let $F = \langle v_1, v_2, \ldots, v_n \rangle$
    \State $\text{startIndex} \gets n - 1$
    \State $\text{endIndex} \gets 0$
    \While{$\text{startIndex} > \text{endIndex} + 1$}
        \If{chord $(v_{\text{startIndex}}, v_{\text{endIndex}})$ exists in $S$}
            \State $\text{startIndex} \gets \text{startIndex} - 1$
        \Else
            \State $\text{endIndex} \gets \text{endIndex} + 1$
        \EndIf
    \EndWhile
    \State \Return $v_{\text{endIndex}}$ \Comment{Safe root vertex found}

\EndProcedure
\end{algorithmic}
\end{algorithm}



\subsection{Algorithm Walkthrough}

To illustrate how the algorithm works, we provide a brief walkthrough:

\begin{enumerate}
    \item \textbf{Initialize:} Start with an empty global chord set $S$ and empty partial triangulation $T$.
    
    \item \textbf{Process Face $F_i$:} 
    \begin{itemize}
        \item Find a safe root vertex $r$ for $F_i$ using \textsc{FindSafeRootVertex}
        \item Construct the root triangulation $T_{ir}$ with all chords from $r$
        \item Add these chords to $S$
    \end{itemize}
    
    \item \textbf{Explore Local Tree:}
    \begin{itemize}
        \item Starting from $T_{ir}$, traverse the local genealogical tree
        \item At each triangulation, either output (if last face) or recurse to next face
        \item Generate children by flipping the generating chord
        \item Maintain $S$ through backtracking (add/remove chords)
    \end{itemize}
    
    \item \textbf{Backtrack:} After exploring all descendants, remove chords from $S$ and try alternative triangulations.
    
    \item \textbf{Complete:} When all faces have been processed with all their triangulation combinations, output the complete triangulation.
\end{enumerate}

\subsection{Key Properties}

The algorithm maintains several important invariants:

\begin{itemize}
    \item \textbf{Global Chord Set Consistency:} At any point, $S$ contains exactly the chords from the current path in the recursion tree (faces $F_1$ through $F_{i-1}$ are fully triangulated, and face $F_i$ has a specific triangulation being explored).
    
    \item \textbf{No Multi-edges:} Before adding any chord, we verify it is not in $S$, ensuring no multi-edges are created.
    
    \item \textbf{Local Tree Structure:} Each face's triangulations form a tree rooted at the safe root triangulation, ensuring systematic exploration without duplicates.
    
    \item \textbf{Completeness:} By exploring all branches in each local tree and all faces, every valid combination is generated exactly once.
\end{itemize}


\section{Correctness and Completeness}
\label{sec:correctness}

We now prove that our algorithm correctly generates all triangulations without duplicates.

\subsection{Existence of Safe Root Vertices}

\begin{theorem}[Existence of Safe Root]
\label{thm:safe_root_exists}
In any face of a biconnected plane graph during the triangulation process, there exist at least two safe root vertices.
\end{theorem}

\begin{proof}
Consider a face $F$ with $n$ vertices $v_1, v_2, \ldots, v_n$ listed in counterclockwise order. Let $S$ be the global chord set from previously processed faces.

\textbf{Case 1:} If $S$ contains no chords with both endpoints in $F$, then all $n$ vertices are safe roots.

\textbf{Case 2:} Assume $S$ contains at least one chord with both endpoints in $F$. Consider a chord $(v_a, v_b) \in S$. Due to planarity, no other chord can cross $(v_a, v_b)$.

The chord $(v_a, v_b)$ divides the cycle $F$ into two segments:
\begin{itemize}
    \item $S_1 = \langle v_a, v_{a+1}, \ldots, v_b \rangle$
    \item $S_2 = \langle v_b, v_{b+1}, \ldots, v_a \rangle$
\end{itemize}

Each segment forms a substructure $T_m$ where $m$ is the number of internal vertices (vertices excluding the endpoints).

\textbf{Base Case: $T_1$}

If $S_1$ has one internal vertex $v_c$ between $v_a$ and $v_b$, then both $v_a$ and $v_b$ are neighbors of $v_c$ on the boundary. Any potential chord from $v_c$ would go to one of its neighbors, which is not allowed. Therefore, $v_c$ has no conflicting chords and is a safe root.

\textbf{Base Case: $T_2$}

If $S_1$ has two internal vertices $v_c$ and $v_d$ (with $v_c$ between $v_a$ and $v_d$, and $v_d$ between $v_c$ and $v_b$), then:
\begin{itemize}
    \item The only non-neighbor of $v_c$ is $v_b$
    \item The only non-neighbor of $v_d$ is $v_a$
\end{itemize}

If $(v_c, v_b) \in S$, this creates a $T_1$ structure containing only $v_d$, making $v_d$ safe.
If $(v_d, v_a) \in S$, this creates a $T_1$ structure containing only $v_c$, making $v_c$ safe.
If neither chord exists, both $v_c$ and $v_d$ are safe.

\textbf{Inductive Case: $T_m$ for $m > 2$}

We use an iterative reduction strategy. Consider the first internal vertex $v_{a+1}$:
\begin{itemize}
    \item Its neighbors on the boundary are $v_a$ and $v_{a+2}$
    \item Possible conflicting chords: $(v_{a+1}, v_{a+3}), \ldots, (v_{a+1}, v_b)$
\end{itemize}

Check from $v_b$ downward:
\begin{itemize}
    \item If $(v_{a+1}, v_b) \in S$, the structure reduces to $T_{m-1}$ with vertices $v_{a+2}, \ldots, v_b$
    \item If $(v_{a+1}, v_{b-k}) \in S$ for some $k > 0$, the structure reduces to $T_{m-k-1}$
    \item If no such chord exists, $v_{a+1}$ is a safe root
\end{itemize}

This process either finds a safe root or reduces $m$, eventually reaching $T_1$ or $T_2$, both of which contain safe roots.

\textbf{Finding Safe Vertices in $O(n)$ Time:}

The above iterative process can be implemented efficiently:

\begin{enumerate}
    \item Consider vertex $v_{a+1}$ (first internal vertex)
    \item $v_{a+1}$ cannot create chords with neighbors $v_a$ and $v_{a+2}$
    \item Possible chord partners: $\{v_{a+3}, \ldots, v_b\}$
    \item Check from $v_b$ downward using a hash set ($O(1)$ per check)
    \item If a chord $(v_{a+1}, v_{b-k})$ is found, reduce to subproblem $T_{m-k-1}$
    \item If no chord is found after checking all candidates, $v_{a+1}$ is safe
\end{enumerate}

Each iteration either:
\begin{itemize}
    \item Finds a safe root vertex (algorithm terminates successfully), or
    \item Reduces the problem size by at least 1, eventually reaching $T_1$ or $T_2$
\end{itemize}

Since $T_1$ and $T_2$ are proven to contain safe vertices, and each reduction step decreases $m$, the algorithm is guaranteed to find a safe vertex in $O(n)$ time.

\textbf{Two Complementary $T_m$ Structures:}

The face $F$ is divided by chord $(v_a, v_b)$ into two segments:
\begin{itemize}
    \item Structure 1: $S_1 = \langle v_a, v_{a+1}, \ldots, v_b \rangle$ with $m_1$ internal vertices
    \item Structure 2: $S_2 = \langle v_b, v_{b+1}, \ldots, v_a \rangle$ with $m_2$ internal vertices
\end{itemize}

Note that $m_1 + m_2 = n - 2$ (total vertices minus the two endpoints of the constraining chord).

By the analysis above, each structure $S_1$ and $S_2$ contains at least one safe vertex. Therefore, every face has at least two safe root vertices—one from each complementary segment.

Since both segments $S_1$ and $S_2$ contain safe roots by this argument, every face has at least two safe root vertices.
\end{proof}

\begin{remark}
The existence of at least two safe root vertices is crucial for the algorithm's flexibility. It ensures that even in highly constrained faces with many pre-existing chords from previous face triangulations, we can always find a valid starting point for the local genealogical tree.
\end{remark}

\subsection{Completeness via Pruning}

\begin{theorem}[Completeness]
\label{thm:completeness}
The algorithm generates all valid triangulations of the biconnected plane graph without duplications.
\end{theorem}

\begin{proof}
We must show two properties:
\begin{enumerate}
    \item Every valid triangulation is generated (completeness)
    \item No triangulation is generated twice (no duplicates)
\end{enumerate}

\textbf{Pruning Preserves Validity:}

When the algorithm prunes a branch (because a chord would create a multi-edge), we must verify that all children of that branch would also be invalid.

Consider a triangulation $T$ with a blocking chord $(v_a, v_b)$ that already exists in the global set $S$. Since $T$ is a complete triangulation of the face, there exist two vertices $v_p$ and $v_q$ forming triangles $(v_a, v_b, v_p)$ and $(v_a, v_b, v_q)$ on opposite sides of the chord.

We now prove that a blocking chord is \emph{never flipped} in any child of the triangulation tree, by analyzing all possible positions of the vertices relative to the root vertex $r$.

\textbf{Case 1: $v_p$ is on the left of the root vertex $r$}

Flipping the $(v_a, v_b)$ chord would result in a $(v_p, v_q)$ chord. However, before this flip, the $(v_b, v_p)$ chord was also a blocking chord positioned to the left of $(v_a, v_b)$ (by the ordering of vertices). This means $(v_a, v_b)$ could not be the leftmost blocking chord. 

Consequently, after flipping to create the new $(v_p, v_q)$ chord, this new chord cannot be the leftmost blocking chord either, because the $(v_p, v_b)$ chord still exists and remains to the left of the $(v_q, v_p)$ chord. 

Therefore, flipping the $(v_a, v_b)$ blocking chord is not permitted by the algorithm in this case, as the algorithm only flips the leftmost generating chord.

\textbf{Case 2: $v_p$ is on the right of the root vertex $r$}

Flipping the $(v_a, v_b)$ chord would produce a $(v_p, v_q)$ chord. However, the $(v_a, v_p)$ chord exists and is positioned to the left of where the new $(v_p, v_q)$ chord would be. Since $(v_a, v_p)$ is more to the left than $(v_p, v_q)$, the flip is invalid according to the algorithm's leftmost-chord-first policy.

\textbf{Case 3: $v_p$ is at the root vertex $r$}

Flipping the $(v_a, v_b)$ chord would create a $(r, v_q)$ chord, where $r$ is the root vertex. This would give the child triangulation's root vertex better visibility than the parent—the root would be able to see $v_q$ through a chord, which contradicts the tree construction where chords from the root are only added in the root triangulation.

This violates the algorithm's fundamental constraint that all chords emanating from the root vertex are established in the root triangulation and are never created through flipping operations in descendant triangulations.

\textbf{Conclusion:}

In all three cases, flipping a blocking chord in a triangulation tree is impossible according to the algorithm's rules. Therefore, if $(v_a, v_b)$ is a blocking chord (duplicate from previous faces), it will be present in all descendants of $T$ in the local tree, making all descendants invalid. 

By pruning such branches, we eliminate no valid triangulations, preserving completeness.

\textbf{No Duplicates:}

The local genealogical tree structure for each face ensures unique parent-child relationships: each triangulation (except the root) has exactly one parent obtained by flipping its generating chord. This defines a tree (not a DAG), preventing cycles and duplicates within a single face's triangulation space.

Across faces, the recursive DFS exploration with backtracking ensures that each combination of face triangulations is explored exactly once. The global chord set $S$ is updated and restored correctly during backtracking, maintaining consistency.

Therefore, the algorithm generates all valid triangulations exactly once.
\end{proof}

\section{Complexity Analysis}
\label{sec:complexity}

\subsection{Time Complexity}

\begin{theorem}[Amortized Constant Time per Triangulation]
\label{thm:time_complexity}
The algorithm has amortized $O(1)$ time complexity per triangulation.
\end{theorem}

\begin{proof}
Let the biconnected graph have $k$ faces with $n_1, n_2, \ldots, n_k$ vertices respectively. Each face $F_i$ has at least $O(n_i)$ triangulations (proven in Lemma~\ref{lem:min_triangulations}).

\textbf{Building Time:}

For face $F_i$ with $n_i$ vertices:
\begin{itemize}
    \item Finding a safe root: $O(n_i)$ time (Algorithm~\ref{alg:safe_root})
    \item Constructing root triangulation: $O(n_i)$ time
    \item Total building time: $B_i = O(n_i)$
\end{itemize}

\textbf{Total Triangulations:}

The total number of complete graph triangulations is at least:
\begin{align}
T \geq \prod_{i=1}^{k} (n_i - 3) = \Omega\left(\prod_{i=1}^{k} n_i\right)
\end{align}

\textbf{Amortized Analysis:}

The total building time across all faces and all recursive calls is:
\begin{align}
B_{\text{total}} &= \sum_{i=1}^{k} B_i \cdot \prod_{j=1}^{i-1} n_j \\
&= B_1 + B_2 \cdot n_1 + B_3 \cdot n_1 \cdot n_2 + \cdots + B_k \cdot n_1 \cdot n_2 \cdots n_{k-1}
\end{align}

\textbf{Two-Face Analysis ($k = 2$):}

For $k = 2$ faces:
\begin{align}
B_{\text{total}} &= n_1 + n_1 \cdot n_2 \\
T &= n_1 \cdot n_2
\end{align}

Since $n_2 \geq 3$ (minimum cycle size):
\begin{align}
n_1 \cdot 1 &< n_1 \cdot n_2 \\
n_1 + n_1 \cdot n_2 &< n_1 \cdot n_2 + n_1 \cdot n_2 \\
B_{\text{total}} &< 2 \cdot n_1 \cdot n_2 = 2T
\end{align}

\textbf{Three-Face Analysis ($k = 3$):}

For $k = 3$ faces with $n_1$, $n_2$, $n_3$ vertices:
\begin{align}
B_{\text{total}} &= n_1 + n_1 \cdot n_2 + n_1 \cdot n_2 \cdot n_3 \\
T &= n_1 \cdot n_2 \cdot n_3
\end{align}

From the two-face case, we know: $n_1 + n_1 \cdot n_2 < 2 \cdot n_1 \cdot n_2$

Therefore:
\begin{align}
B_{\text{total}} &< 2 \cdot n_1 \cdot n_2 + n_1 \cdot n_2 \cdot n_3
\end{align}

Since $n_3 \geq 3$, we have $2 < n_3$:
\begin{align}
2 \cdot n_1 \cdot n_2 &< n_1 \cdot n_2 \cdot n_3 \\
B_{\text{total}} &< n_1 \cdot n_2 \cdot n_3 + n_1 \cdot n_2 \cdot n_3 \\
&= 2T
\end{align}

\textbf{General Case (Induction):}

\begin{lemma}[Building Time Bound]
For any number of faces $k \geq 2$ where each face has at least 3 vertices:
\begin{align}
B_{\text{total}} < 2T
\end{align}
\end{lemma}

\begin{proof}[Proof by Induction]
\emph{Base case:} $k = 2$ is proven above.

\emph{Inductive hypothesis:} Assume for $k - 1$ faces:
\begin{align}
n_1 + n_1 \cdot n_2 + \cdots + n_1 \cdot n_2 \cdots n_{k-1} < 2 \cdot n_1 \cdot n_2 \cdots n_{k-1}
\end{align}

\emph{Inductive step:} For $k$ faces:
\begin{align}
B_{\text{total}} &= n_1 + n_1 \cdot n_2 + \cdots + n_1 \cdot n_2 \cdots n_{k-1} + n_1 \cdot n_2 \cdots n_k \\
&< 2 \cdot n_1 \cdot n_2 \cdots n_{k-1} + n_1 \cdot n_2 \cdots n_k \quad \text{(by hypothesis)}
\end{align}

Since $n_k \geq 3$:
\begin{align}
2 \cdot n_1 \cdot n_2 \cdots n_{k-1} &< n_1 \cdot n_2 \cdots n_{k-1} \cdot n_k \\
B_{\text{total}} &< n_1 \cdot n_2 \cdots n_k + n_1 \cdot n_2 \cdots n_k \\
&= 2 \cdot n_1 \cdot n_2 \cdots n_k = 2T
\end{align}
\end{proof}

Including output time (constant per triangulation), the total time is $< 3T$, yielding amortized $O(1)$ time per triangulation.
\end{proof}

\begin{corollary}[Amortized Linear Time]
For any biconnected planar graph with $k$ faces, the algorithm takes total time $< 3 \times (\text{number of triangulations})$. Therefore, the algorithm has amortized linear time complexity: $O(1)$ per triangulation, or $O(T)$ overall where $T$ is the total number of triangulations.
\end{corollary}

\begin{lemma}[Minimum Triangulation Count]
\label{lem:min_triangulations}
A face with $n$ vertices has at least $n - 3$ distinct triangulations.
\end{lemma}

\begin{proof}
By Theorem~\ref{thm:safe_root_exists}, every face has at least two safe root vertices $v_a$ and $v_c$. The $v_a$-root triangulation has all chords emanating from $v_a$, while the $v_c$-root triangulation has all chords from $v_c$.

At most one chord $(v_a, v_c)$ can appear in both triangulations. All other $n - 4$ chords differ. Transforming one into the other via edge flips requires at least $n - 4$ flips, each producing a distinct valid triangulation.

Including the root: at least $n - 3$ triangulations exist.
\end{proof}

\subsection{Space Complexity}

\begin{theorem}[Linear Space Complexity]
\label{thm:space_complexity}
The algorithm uses $O(n)$ space, where $n$ is the total number of vertices.
\end{theorem}

\begin{proof}
The algorithm stores:
\begin{itemize}
    \item Vertex information for all faces: $O(n)$
    \item Global chord set $S$: A planar graph has at most $3n - 6$ edges, so $O(n)$ chords
    \item Recursion stack depth: $O(k)$ where $k$ is the number of faces, and $k \leq n$
    \item Current triangulation chords: $O(n)$
\end{itemize}

Total space: $O(n)$.
\end{proof}

\section{Implementation and Experimental Results}
\label{sec:implementation}

We implemented the algorithm in Python, with comprehensive testing on various biconnected plane graphs. The implementation is available at \url{https://github.com/TAR2003/ThesisGraph}.

\subsection{Implementation Details}

Key implementation aspects include:
\begin{itemize}
    \item \textbf{Face Representation:} Each face is stored as a list of vertices in counterclockwise order
    \item \textbf{Global Chord Set:} Implemented using a hash set for $O(1)$ lookup and insertion
    \item \textbf{Triangulation Representation:} Each triangulation is a set of chords (vertex pairs)
    \item \textbf{Edge Flipping:} Efficiently computed by identifying the generating chord and its adjacent triangles
    \item \textbf{Canonical Form:} Each triangulation is stored in a canonical form (sorted chord list) to facilitate duplicate detection and verification
\end{itemize}

\subsection{DFS Exploration Tree Structure}

The algorithm explores triangulations using a depth-first search (DFS) tree where:

\begin{itemize}
    \item \textbf{Level 0:} Root node (no faces processed yet)
    \item \textbf{Level $i$ ($1 \leq i \leq k$):} Nodes represent partial triangulations where faces $F_1, \ldots, F_i$ are triangulated
    \item \textbf{Level $k$:} Leaf nodes represent complete triangulations (all faces processed)
\end{itemize}

At each level $i$, the algorithm:
\begin{enumerate}
    \item Finds a safe root vertex for face $F_i$
    \item Constructs the root triangulation
    \item Explores the local genealogical tree of face $F_i$ via DFS
    \item For each triangulation of $F_i$, recursively processes face $F_{i+1}$
    \item Backtracks by removing chords from the global set
\end{enumerate}

This creates a multi-level tree where each path from root to leaf represents one complete triangulation of the entire graph.

\subsection{GlobalChordSet Management}

The GlobalChordSet $S$ plays a crucial role in preventing multi-edge creation:

\begin{itemize}
    \item \textbf{Initialization:} $S = \emptyset$ before processing any face
    \item \textbf{Addition:} When exploring a triangulation $T$ of face $F_i$, add all chords of $T$ to $S$
    \item \textbf{Checking:} Before adding any chord, verify it does not already exist in $S$
    \item \textbf{Removal:} After processing all descendants (backtracking), remove chords of $T$ from $S$
    \item \textbf{Restoration:} This ensures correct state for exploring alternative triangulations
\end{itemize}

The dynamic addition and removal of chords during DFS ensures that $S$ always contains exactly the chords from the current path in the exploration tree, enabling efficient conflict detection.

\subsection{Test Cases}

We tested the algorithm on various graph structures:
\begin{itemize}
    \item \textbf{Simple polygons:} Triangle, square, pentagon, hexagon, heptagon, octagon, nonagon
    \item \textbf{Biconnected graphs with 2 faces:} Various configurations including bow-tie, double-square patterns
    \item \textbf{Complex biconnected structures:} Graphs with 3+ faces sharing edges
    \item \textbf{Graphs with varying connectivity:} From minimal biconnected to near-triconnected graphs
    \item \textbf{Degenerate cases:} Graphs with many pre-existing chords from the input
    \item \textbf{Stress tests:} Large graphs with dozens of vertices and multiple faces
\end{itemize}

Example test inputs from our test suite:
\begin{itemize}
    \item \textbf{Hexagon:} Single 6-vertex cycle, expected: 14 triangulations
    \item \textbf{Bow-tie:} Two squares sharing a vertex, 2 faces
    \item \textbf{Double-square:} Two squares sharing an edge, 2 faces
    \item \textbf{Complex wheel:} Hub vertex connected to outer cycle with 5 spokes
\end{itemize}

\subsection{Performance Results}

Experimental results confirm the theoretical analysis:
\begin{itemize}
    \item \textbf{Completeness:} All triangulations were generated without duplicates (verified by exhaustive enumeration for small instances and canonical form checking for larger instances)
    \item \textbf{Efficiency:} Average time per triangulation remained constant as graph size increased, confirming $O(1)$ amortized time
    \item \textbf{Scalability:} The algorithm successfully handled graphs with dozens of vertices and multiple faces, generating thousands of triangulations efficiently
    \item \textbf{Memory usage:} Linear space complexity confirmed—memory usage scaled linearly with the number of vertices
    \item \textbf{Safe root finding:} Always found safe root vertices in $O(n)$ time, with most faces having multiple safe root options
\end{itemize}

\subsection{Verification and Validation}

To ensure correctness, we implemented several validation mechanisms:

\begin{enumerate}
    \item \textbf{Triangulation Validity:} Each output is verified to be a valid triangulation (all faces are triangles, no crossing chords)
    \item \textbf{Duplicate Detection:} Hash-based duplicate detection confirms no triangulation is output twice
    \item \textbf{Completeness Checking:} For small graphs, we compare output count with known theoretical bounds
    \item \textbf{Multi-edge Detection:} Verify that no multi-edges exist in any output triangulation
    \item \textbf{Planarity Verification:} Confirm all generated triangulations maintain planarity
\end{enumerate}

All tests passed successfully, confirming the algorithm's correctness and efficiency.

\section{Conclusion and Future Work}
\label{sec:conclusion}

We have presented the first algorithm for generating all triangulations of biconnected plane graphs with constant amortized time complexity. Our algorithm extends the elegant tree-of-triangulations framework of Parvez, Rahman, and Nakano~\cite{parvez2011} from triconnected to biconnected graphs by introducing safe root selection and local genealogical trees.

\subsection{Key Achievements}

\begin{enumerate}
\item \textbf{Theoretical Foundation:} We have provided rigorous proofs establishing:
    \begin{itemize}
        \item Existence of at least two safe root vertices in every face (Theorem~\ref{thm:safe_root_exists})
        \item Completeness of the algorithm without duplications (Theorem~\ref{thm:completeness})
        \item Amortized $O(1)$ time complexity per triangulation (Theorem~\ref{thm:time_complexity})
        \item Linear $O(n)$ space complexity (Theorem~\ref{thm:space_complexity})
    \end{itemize}

\item \textbf{Practical Algorithm:} Efficient implementation suitable for real-world applications with extensive experimental validation

\item \textbf{Novel Techniques:} Safe root selection and global chord management strategies that can extend to other enumeration problems involving planar graph decompositions

\item \textbf{Broader Applicability:} Extension from triconnected to the more general class of biconnected plane graphs, significantly expanding the scope of graphs that can be processed efficiently
\end{enumerate}

\subsection{Theoretical Significance}

This work makes several theoretical contributions:

\begin{itemize}
    \item \textbf{Generalization:} We show that constant-time triangulation enumeration is achievable beyond the restrictive triconnected case
    \item \textbf{Multi-edge Resolution:} Our safe root concept provides a general strategy for handling chord conflicts in face-decomposed planar graphs
    \item \textbf{Amortized Analysis:} The detailed amortized analysis demonstrates that building time is dominated by output size even with complex face interactions
    \item \textbf{Structural Insights:} The $T_m$ structure analysis reveals fundamental properties of planar graph faces under chord constraints
\end{itemize}

\subsection{Future Directions}

Several promising research directions remain:

\begin{enumerate}
\item \textbf{General Planar Graphs:} Extend the algorithm to simply connected (but not biconnected) planar graphs by handling cut vertices. This would require decomposing the graph at cut vertices and combining triangulations across components.

\item \textbf{Constrained Triangulations:} Generate triangulations satisfying additional constraints such as:
    \begin{itemize}
        \item Minimum angle constraints (for quality mesh generation)
        \item Delaunay property (empty circle criterion)
        \item Weight minimization (for optimal triangulations)
        \item Forbidden edge constraints
    \end{itemize}

\item \textbf{Parallel Enumeration:} Exploit the tree structure for parallel generation of triangulations. The local genealogical trees for different faces can potentially be explored in parallel, with synchronization on the global chord set.

\item \textbf{Applications:} Apply the algorithm to:
    \begin{itemize}
        \item Mesh generation for finite element analysis
        \item VLSI layout and floorplanning optimization
        \item Graph visualization and aesthetic drawing
        \item Computational origami and folding patterns
        \item Computational biology (protein structure analysis)
    \end{itemize}

\item \textbf{Unlabeled Triangulations:} Extend to generate unlabeled triangulations modulo graph isomorphism, reducing redundancy for symmetric graphs

\item \textbf{Counting Triangulations:} Develop efficient algorithms for \emph{counting} (without enumerating) all triangulations, potentially using dynamic programming on the face structure

\item \textbf{Random Sampling:} Design algorithms for uniformly random sampling from the space of all triangulations, useful for probabilistic analysis and Monte Carlo methods

\item \textbf{Dynamic Triangulation:} Investigate online algorithms that maintain the set of triangulations under graph modifications (edge/vertex additions/deletions)
\end{enumerate}

\subsection{Impact}

This work contributes to the theoretical understanding of planar graph triangulations and provides a practical tool for computational geometry applications. By solving the multi-edge problem inherent in biconnected graphs, we have broadened the applicability of constant-time triangulation enumeration beyond the restrictive triconnected case.

The safe root vertex concept and local genealogical tree structure introduced here may find applications in other graph enumeration problems where global constraints must be maintained across local decompositions. The techniques developed for managing the global chord set during recursive exploration are applicable to any decomposition-based enumeration algorithm.

\subsection{Open Problems}

We conclude by highlighting several open questions:

\begin{enumerate}
    \item What is the exact number of safe root vertices in a face as a function of its size and the number of constraining chords?
    \item Can the safe root finding algorithm be improved to $o(n)$ time through preprocessing or advanced data structures?
    \item Is there a closed-form formula for the number of triangulations of a biconnected graph given its face structure?
    \item Can the algorithm be adapted to non-planar graphs with a fixed embedding on other surfaces (torus, projective plane)?
\end{enumerate}

\section*{Acknowledgments}

The author gratefully acknowledges the foundational work of Mohammad Tanvir Parvez, Md. Saidur Rahman, and Shin-ichi Nakano, whose elegant algorithm for triconnected plane graphs inspired this research. Special thanks to the faculty and students at the Department of Computer Science and Engineering, Bangladesh University of Engineering and Technology (BUET), for their valuable feedback and support.

\bibliographystyle{plain}
\bibliography{references}

\end{document}
